\begin{figure}[H]\centering
\begin{tikzpicture}[scale=0.62]
\othbg{6}{6}
% pozycja startowa 6x6 (czarny rusza): biały na przekątnej głównej, czarny na przeciwprzekątnej
\othW{2.5}{3.5}\othB{3.5}{3.5}
\othB{2.5}{2.5}\othW{3.5}{2.5}
% etykiety kolumn (a--f) i wierszy (1--6, od góry)
\foreach \c/\l in {0/a,1/b,2/c,3/d,4/e,5/f} \node[below,font=\scriptsize] at (\c+0.5,-0.05) {\l};
\foreach \y/\l in {5.5/1,4.5/2,3.5/3,2.5/4,1.5/5,0.5/6} \node[left,font=\scriptsize] at (-0.05,\y) {\l};
\end{tikzpicture}
\caption{Pozycja startowa Othello 6$\times$6: cztery pionki w~układzie krzyżowym w~centrum
(biały na przekątnej głównej, czarny na przeciwprzekątnej). Czarny wykonuje pierwszy ruch.}
\label{fig:board}
\end{figure}
